\documentclass{article}
\usepackage[T1]{fontenc}
\usepackage{lmodern}
\usepackage{graphicx}
\usepackage{amsmath,amssymb}
\usepackage[a4paper,margin=2cm]{geometry}
\usepackage[absolute,overlay]{textpos}
\setlength{\TPHorizModule}{1mm}
\setlength{\TPVertModule}{1mm}
\usepackage[indonesian]{babel}
\usepackage{hyperref}
\setlength{\emergencystretch}{2em}
% unit: Halaman 27

\begin{document}

% === Halaman 27 (llm) ===

\section*{B. Vig\`enere Cipher}

\begin{itemize}
    \item Vigenere Cipher termasuk ke dalam kelompok \textit{cipher} abjad-majemuk (\textit{polyalphabetic cipher}). Ini berbeda dengan cipher abjad-tunggal (\textit{monolphabetic cipher}).
    \item Perbedaan \textit{polyalphabetic cipher} dengan \textit{monoalphabetic cipher}:
    \begin{itemize}
        \item[$\blacktriangleright$] \textit{Cipher} abjad-tunggal: satu kunci untuk semua huruf plainteks \\ Contoh: Caesar cipher
        \item[$\blacktriangleright$] \textit{Cipher} abjad-majemuk: setiap huruf menggunakan kunci berbeda. \\ Contoh: Vigenere cipher
    \end{itemize}
    \item \textit{Cipher} abjad-majemuk dibuat untuk mengatasi kelemahan \textit{cipher} abjad-Tunggal yang mudah dipecahkan dengan teknik analisis frekuensi
\end{itemize}

\end{document}
